\section{Discussion}
\label{sec:discussion}

The economic model presented in Sections~\ref{sec:tco}--\ref{sec:credits}
reveals a nuanced adoption landscape for verifiable AI accountability.
While the marginal cost of transparency approaches zero at scale, the
fixed-cost structure creates distinct strategic dynamics worth examining
beyond the crossover arithmetic.

\subsection{The Valley of Non-Compliance}
\label{subsec:valley}

Our analysis identifies a structural gap for small-to-medium enterprises
(SMEs). For operators processing fewer than $\sim$250{,}000
decisions/year, the fixed infrastructure cost (\$5{,}000) exceeds
the expected penalty savings under the compliance-credit model at any
regime except SEC-regulated finance. This creates a \emph{Valley of
Non-Compliance}: a range of deployment scales where rational economic
actors are incentivised to remain opaque, relying on low enforcement
probability as de-facto protection.

This finding contradicts the implicit assumption in the EU AI Act that
compliance is uniformly accessible. Without targeted intervention ---
shared infrastructure, subsidised log servers, or a SME compliance-credit
transfer market (Section~\ref{sec:credits}) --- the BTV stack risks
becoming a competitive moat available only to large enterprises,
reinforcing market concentration rather than diffusing accountability.

\subsection{Adoption Barriers Beyond Cost}
\label{subsec:barriers}

Even when the economic case is favourable (for $N > N^*_{\text{credit}}$),
non-economic barriers persist:

\begin{itemize}
  \item \textbf{Technical friction.} The BTV stack requires Safe Rust
  and systems-level engineering (Papers~1--2~\cite{btv2026,btv-t-repo}).
  The prevailing AI development paradigm --- Python-based ML pipelines
  without ownership semantics --- lacks the compile-time type structure
  required by the stack. Adoption requires an engineering culture shift
  analogous to the managed/unmanaged code transition of the early 2000s.
  The interview data (Section~\ref{sec:interviews}) will quantify how
  frequently compliance officers identify this as the primary barrier.

  \item \textbf{Strategic opacity.} Opacity is not merely a cost-saving
  measure; for some operators it is a defence mechanism. Verifiable
  audit trails expose algorithmic bias and strategic errors that would
  otherwise remain within plausible-deniability range. Internal
  resistance from product and legal teams is likely independent of
  the economic calculus.
\end{itemize}

\subsection{Sensitivity of the Crossover to Enforcement Credibility}
\label{subsec:sensitivity}

The crossover point $N^*_{\text{credit}}$ is inversely proportional to
the risk density $\rho$, which in turn depends on enforcement
credibility. A 50\% reduction in enforcement probability --- plausible
if AI Act enforcement ramps slowly or if regulatory capture reduces
fine magnitudes --- doubles the crossover threshold from 500{,}000 to
1{,}000{,}000 decisions/year, pricing out a significant fraction of
mid-market AI operators.

This creates a dependence structure between technical and institutional
choices that the compliance-credit mechanism partially resolves:
operators who invest early lock in safe-harbour protection regardless
of subsequent enforcement trajectory, reducing their exposure to
regulatory uncertainty. In this sense, early adoption functions as
a hedge against enforcement volatility, not only against enforcement
itself.

\subsection{Limitations}
\label{subsec:limitations}

\begin{enumerate}
  \item \textbf{TCO model.} Fixed infrastructure costs are estimated
  from managed cloud pricing as of 2026 and may not reflect
  on-premises deployments or sovereign-cloud constraints in regulated
  jurisdictions.

  \item \textbf{Enforcement probability.} $P_{\text{enf}}$ estimates
  are derived from historical averages; AI Act enforcement trajectory
  is speculative, and sample sizes for Art.~22 GDPR actions remain
  small.

  \item \textbf{Interview sample.} $N=12$ is a small, purposive sample.
  Results from Section~\ref{sec:interviews} should be treated as
  hypothesis-generating, not representative.

  \item \textbf{Linear cost scaling.} The model assumes linear scaling
  of variable costs. Log infrastructure may exhibit non-linear cost at
  $10^8$+ scale due to sharding and replication requirements not
  captured in Table~\ref{tab:tco}.

  \item \textbf{$\delta$ estimate.} The 50\% mitigation factor is
  grounded in existing enforcement decisions but has not been formally
  validated with supervisory authorities. Regulatory recognition of
  cryptographic compliance guarantees as ``appropriate technical
  measures'' under GDPR Art.~25 remains an open interpretive question.
\end{enumerate}
